El presente capítulo introduce el marco general en el que se inscribe este trabajo final, estableciendo los fundamentos conceptuales, tecnológicos y profesionales que justifican su desarrollo. En un contexto marcado por una acelerada evolución de la inteligencia artificial y, en particular, por la irrupción de los \glspl{llm}, se vuelve necesario analizar de qué manera estas tecnologías emergentes pueden ser aplicadas de forma significativa y responsable en ámbitos específicos de formación y capacitación profesional \parencite{zhao2023survey, gallegos2024bias}.

Asimismo, este capítulo aborda la relevancia de la certificación \gls{pmp} como estándar internacional en gestión de proyectos, destacando tanto su valor en el mercado laboral como las exigencias asociadas a su preparación. A partir de esta convergencia entre tecnología y formación profesional, se plantea el problema central que motiva el trabajo: las limitaciones de los enfoques tradicionales de estudio y la oportunidad de incorporar asistentes inteligentes basados en \gls{llm} como apoyo al aprendizaje.

En este marco, se formulan la hipótesis principal y las preguntas de investigación que guían el diseño, desarrollo y evaluación del asistente virtual propuesto, delimitando con claridad el alcance del trabajo y diferenciando aquello que forma parte del estudio de aquello que queda explícitamente excluido. Finalmente, se presenta la organización general del documento, ofreciendo al lector una visión estructurada del contenido de los capítulos siguientes y facilitando la comprensión del recorrido metodológico y conceptual de la investigación.

\section{Contexto y motivación del trabajo}

El presente trabajo se inscribe en un contexto de profundas transformaciones tecnológicas, organizacionales y educativas impulsadas por el avance acelerado de la \gls{ia}, y en particular por la irrupción de los \glspl{llm}. Estas tecnologías han comenzado a modificar de manera sustancial la forma en que las personas acceden al conocimiento, aprenden nuevas competencias y se preparan para desafíos profesionales cada vez más complejos \parencite{zhao2023survey}.

En paralelo, el mundo profesional demanda perfiles capaces de gestionar proyectos con altos niveles de complejidad, incertidumbre y multidisciplinariedad. En este escenario, las certificaciones internacionales, como la \gls{pmp}, se han consolidado como estándares de referencia para validar competencias en gestión de proyectos. Sin embargo, la preparación para este tipo de certificaciones continúa siendo un proceso exigente, intensivo y, en muchos casos, poco personalizado.

La convergencia entre estos dos fenómenos —el avance de los \glspl{llm} y la creciente demanda de formación profesional especializada— configura el marco que motiva el desarrollo del presente trabajo. En particular, se identifica una oportunidad clara para explorar el uso de asistentes virtuales inteligentes como soporte educativo en la preparación para el examen de certificación \gls{pmp}, integrando capacidades avanzadas de procesamiento del lenguaje natural con enfoques pedagógicos orientados al aprendizaje autónomo y adaptativo.

\subsection{Evolución de la inteligencia artificial y los \glspl{llm}}

La inteligencia artificial ha experimentado una evolución sostenida desde sus primeras formulaciones teóricas hasta convertirse en una tecnología de uso transversal en múltiples dominios. En sus etapas iniciales, la \gls{ia} estuvo dominada por enfoques simbólicos y sistemas basados en reglas explícitas, los cuales presentaban importantes limitaciones en términos de escalabilidad, flexibilidad y adaptación a contextos no estructurados.

Con el surgimiento del aprendizaje automático y, posteriormente, del aprendizaje profundo (\textit{deep learning}), la \gls{ia} dio un salto cualitativo significativo. En este marco, los modelos neuronales comenzaron a destacarse por su capacidad para aprender patrones complejos a partir de grandes volúmenes de datos. Esta evolución alcanzó un punto de inflexión con la aparición de los \glspl{llm}, entrenados sobre corpus masivos de texto y capaces de capturar relaciones semánticas, contextuales y pragmáticas del lenguaje natural \parencite{zhao2023survey}.

Modelos como GPT-3 y sus sucesores demostraron que es posible abordar tareas de comprensión, generación y razonamiento en lenguaje natural con niveles de desempeño que, en ciertos contextos, resultan comparables a los humanos \parencite{zhao2023survey}. Posteriormente, la incorporación de capacidades multimodales amplió aún más el alcance de estas tecnologías, habilitando aplicaciones más ricas y sofisticadas.

Los \glspl{llm} se caracterizan no solo por su tamaño y capacidad de generalización, sino también por su habilidad para adaptarse a distintos dominios mediante técnicas como el \textit{prompting} y el \textit{few-shot learning}. Estas propiedades los convierten en una base tecnológica especialmente atractiva para el desarrollo de sistemas interactivos orientados al soporte cognitivo, la asistencia educativa y el acompañamiento en procesos de aprendizaje complejos \parencite{zhao2023survey}.

\subsection{Impacto de los \glspl{llm} en educación y formación profesional}

El ámbito educativo ha sido uno de los primeros en explorar activamente el potencial de los \glspl{llm}. Diversos estudios recientes destacan su capacidad para actuar como tutores virtuales, generar explicaciones personalizadas y ofrecer retroalimentación inmediata, favoreciendo enfoques de aprendizaje activo y reflexivo \parencite{kasneci2023chatgpt}.

En el contexto de la formación profesional, los \glspl{llm} ofrecen ventajas adicionales, como la posibilidad de contextualizar ejemplos en escenarios reales, adaptar las respuestas al nivel de conocimiento del usuario y simular situaciones prácticas complejas. Estas características resultan particularmente relevantes en dominios que requieren la aplicación de criterios situacionales y toma de decisiones fundamentadas.

No obstante, la incorporación de \glspl{llm} en procesos educativos también plantea desafíos relevantes. La literatura coincide en que el mayor valor de estos sistemas no reside en reemplazar al docente ni en sustituir el esfuerzo cognitivo del estudiante, sino en complementar los procesos educativos mediante enfoques pedagógicos bien definidos y controlados \parencite{kasneci2023chatgpt}.

Dentro de este marco, el uso de \gls{llm} como soporte para la preparación de exámenes de certificación profesional constituye un área de investigación aún emergente, con un potencial significativo para mejorar la comprensión conceptual y el razonamiento aplicado.

\subsection{Importancia de la certificación \gls{pmp} en el ámbito profesional}

La certificación \gls{pmp}, otorgada por el \gls{pmi}, es reconocida a nivel internacional como una de las credenciales más prestigiosas en el campo de la gestión de proyectos \parencite{pmi2021handbook}. Su obtención acredita que el profesional posee conocimientos sólidos, experiencia práctica y competencias alineadas con estándares globales de dirección de proyectos.

El examen \gls{pmp} evalúa un amplio espectro de dominios vinculados con personas, procesos y entornos de negocio, integrando enfoques predictivos, ágiles e híbridos, en coherencia con los principios y estándares promovidos por el \gls{pmi} \parencite{pmbok7}. Esta amplitud conceptual, sumada a la profundidad requerida en cada área, convierte a la preparación para el examen en un desafío significativo incluso para profesionales con experiencia previa.

Desde una perspectiva laboral, contar con la certificación \gls{pmp} suele asociarse a mayores oportunidades de empleo, mejores niveles salariales y una mayor credibilidad profesional. Sin embargo, el proceso de preparación continúa apoyándose mayormente en materiales extensos y métodos de estudio poco personalizados, lo que refuerza la necesidad de explorar enfoques alternativos y complementarios.

En este contexto, surge la motivación central del presente trabajo: investigar y desarrollar un asistente virtual basado en \gls{llm} que actúe como soporte inteligente para la preparación del examen \gls{pmp}, ofreciendo explicaciones contextualizadas, retroalimentación personalizada y acompañamiento continuo.

\section{Planteamiento del problema}

La preparación para el examen de certificación \gls{pmp} representa un desafío significativo para una amplia variedad de aspirantes, incluso para aquellos que cuentan con experiencia profesional previa en gestión de proyectos. Si bien el contenido del examen se fundamenta en estándares ampliamente difundidos, como la Guía del \gls{pmbok} y otros marcos desarrollados por el \gls{pmi}, la forma en que dicho contenido es evaluado introduce complejidades adicionales que no siempre son abordadas adecuadamente por los métodos tradicionales de estudio \parencite{pmbok7, pmi2021handbook}.

El examen \gls{pmp} se caracteriza por un enfoque predominantemente situacional, en el cual se presentan escenarios complejos y se solicita al aspirante seleccionar la mejor alternativa de acción entre varias opciones plausibles. Este tipo de evaluación requiere no solo conocimientos teóricos, sino también la capacidad de interpretar contextos, identificar prioridades y aplicar criterios alineados con los principios del \gls{pmi}, más allá de la experiencia personal o de enfoques organizacionales particulares \parencite{pmi2021handbook}.

A partir de estas observaciones, se identifica un problema central vinculado a la brecha existente entre los enfoques tradicionales de preparación y las exigencias cognitivas y conceptuales del examen \gls{pmp}. Esta brecha se manifiesta en dificultades recurrentes para comprender el criterio subyacente a las respuestas correctas, interpretar adecuadamente los escenarios planteados y desarrollar un marco mental consistente con el enfoque evaluativo del \gls{pmi}.

\subsection{Dificultades actuales en la preparación del examen \gls{pmp}}

Entre las principales dificultades que enfrentan los aspirantes al examen \gls{pmp} se encuentra la interpretación precisa de las preguntas situacionales. En muchos casos, los escenarios presentados incluyen información ambigua o incompleta, lo que exige al candidato realizar inferencias, identificar supuestos implícitos y priorizar determinados aspectos por sobre otros.

Asimismo, las opciones de respuesta suelen estar diseñadas de manera que varias de ellas resulten razonables desde distintos enfoques de gestión, lo que incrementa la complejidad del proceso de selección. Esta característica del examen genera confusión en los aspirantes, quienes a menudo encuentran dificultades para identificar cuál de las alternativas se alinea de forma más adecuada con los principios, valores y dominios definidos por el \gls{pmi} \parencite{pmi2021handbook}.

Otra dificultad relevante se relaciona con la comprensión del razonamiento detrás de las respuestas correctas. Muchos materiales de estudio se limitan a indicar la opción correcta sin desarrollar de manera suficiente el fundamento conceptual que justifica dicha elección. Como resultado, el aspirante puede memorizar respuestas sin internalizar el criterio de decisión esperado, reduciendo la efectividad del proceso de aprendizaje y dificultando la transferencia de conocimientos a nuevas situaciones.

\subsection{Limitaciones de los métodos tradicionales de estudio}

Los métodos tradicionales de preparación para el examen \gls{pmp} se apoyan, en gran medida, en libros especializados, cursos presenciales o virtuales, videos grabados y bancos de preguntas estandarizados. Si bien estos recursos han demostrado ser útiles para numerosos candidatos, presentan limitaciones relevantes cuando se analizan desde una perspectiva pedagógica y tecnológica contemporánea.

En primer lugar, estos métodos suelen ofrecer contenidos estáticos, diseñados para un público amplio y heterogéneo, sin contemplar de manera sistemática las diferencias individuales en estilos de aprendizaje, conocimientos previos o áreas específicas de dificultad. Como consecuencia, el aspirante debe realizar un esfuerzo adicional para identificar qué contenidos reforzar y cómo hacerlo, lo que puede resultar ineficiente o desmotivador.

En segundo lugar, la mayoría de los recursos tradicionales presenta un nivel limitado de interactividad. La interacción se reduce, en general, a la lectura pasiva o a la resolución de preguntas con respuestas predefinidas, sin posibilidad de profundizar en un concepto mediante un diálogo dinámico o de explorar distintas líneas de razonamiento en función de las dudas emergentes del estudiante.

Además, estos enfoques carecen de mecanismos de adaptación en tiempo real. No ajustan automáticamente la dificultad de los contenidos, el tipo de ejemplos utilizados ni la profundidad de las explicaciones según el progreso del aspirante, lo que puede generar tanto sobrecarga cognitiva como falta de estímulo, dependiendo del perfil del usuario.

Finalmente, los métodos tradicionales tienden a reproducir modelos de aprendizaje lineales y secuenciales que no siempre se ajustan a las necesidades de los profesionales adultos, quienes suelen beneficiarse de enfoques más flexibles, iterativos y orientados a la resolución de problemas concretos.

\subsection{Oportunidad de uso de \gls{llm} como soporte educativo inteligente}

En este contexto, los \glspl{llm} emergen como una oportunidad tecnológica relevante para abordar las limitaciones identificadas en los métodos tradicionales de preparación para el examen \gls{pmp}. Estos modelos poseen la capacidad de comprender lenguaje natural, analizar contextos complejos y generar explicaciones coherentes, lo que los convierte en una base adecuada para el desarrollo de asistentes educativos inteligentes \parencite{zhao2023survey}.

El uso de LLM como soporte educativo permite concebir entornos de aprendizaje interactivos y adaptativos, en los cuales el aspirante puede formular preguntas en lenguaje natural, solicitar aclaraciones adicionales, explorar ejemplos contextualizados y recibir retroalimentación inmediata ajustada a su nivel de comprensión. Este tipo de interacción favorece un aprendizaje más profundo y reflexivo, especialmente en dominios conceptuales complejos como la gestión de proyectos.

Asimismo, los LLM posibilitan la personalización dinámica del proceso de estudio, al identificar patrones en las consultas del usuario y detectar áreas de dificultad recurrentes. A partir de esta información, el asistente puede adaptar sus explicaciones, proponer nuevos ejemplos o profundizar en determinados conceptos, contribuyendo a optimizar el tiempo de estudio y mejorar la eficiencia del aprendizaje.

Desde una perspectiva pedagógica, la incorporación de \gls{llm} habilita enfoques de aprendizaje activo, tales como el razonamiento guiado, la simulación de escenarios y la reflexión metacognitiva. Estas estrategias resultan particularmente valiosas en el contexto del examen \gls{pmp}, donde comprender el razonamiento detrás de cada decisión es tan importante como seleccionar la respuesta correcta \parencite{kasneci2023chatgpt}.

No obstante, el uso de \gls{llm} en contextos educativos también plantea desafíos relacionados con la confiabilidad de la información, la alineación con los estándares oficiales del \gls{pmi} y la necesidad de promover un uso responsable de la tecnología. Por este motivo, resulta fundamental diseñar el asistente virtual con criterios pedagógicos explícitos, estableciendo límites claros a su funcionamiento y definiendo su rol como complemento —y no sustituto— de los recursos tradicionales de estudio.

\section{Hipótesis y preguntas de investigación}

El presente trabajo se apoya en la premisa de que la inteligencia artificial generativa, y en particular los \glspl{llm}, pueden aportar valor significativo a los procesos de preparación para exámenes de certificación profesional de alta exigencia, como el examen \gls{pmp}, siempre que su uso se diseñe con criterios pedagógicos explícitos y un alcance claramente delimitado.

A partir del análisis del contexto, del problema identificado y de las oportunidades tecnológicas descritas en las secciones anteriores, se formula una hipótesis central que orienta el diseño, desarrollo y evaluación del asistente virtual propuesto. Dicha hipótesis se complementa con un conjunto de preguntas de investigación que permiten abordar el problema desde una perspectiva exploratoria y aplicada.

\subsection{Hipótesis principal}

La hipótesis principal de este trabajo sostiene que un asistente virtual basado en \glspl{llm}, diseñado con un enfoque pedagógico orientado al razonamiento situacional, puede mejorar la comprensión de los escenarios evaluados en el examen \gls{pmp} y facilitar la internalización del criterio de decisión promovido por el \gls{pmi}.

Esta hipótesis no plantea que el asistente garantice la aprobación del examen ni que sustituya los métodos tradicionales de estudio, sino que actúe como un complemento que potencie el proceso de aprendizaje, especialmente en aquellos aspectos relacionados con la interpretación de situaciones complejas, la evaluación de alternativas plausibles y la adopción de un marco mental alineado con los principios del \gls{pmi}.

\subsection{Preguntas de investigación}

Con el objetivo de analizar la validez de la hipótesis planteada y de orientar el desarrollo del asistente virtual, se definen las siguientes preguntas de investigación:

\begin{itemize}
  \item ¿En qué medida un asistente virtual basado en \gls{llm} puede apoyar la comprensión de preguntas situacionales del examen \gls{pmp}?
  \item ¿Cómo contribuyen las explicaciones generadas por el asistente a la internalización del criterio de decisión promovido por el \gls{pmi}?
  \item ¿Cuál es la percepción de los usuarios respecto de la utilidad del asistente como complemento a los métodos tradicionales de preparación para el examen \gls{pmp}?
  \item ¿Cuáles son las principales limitaciones observadas en el uso de un asistente basado en \gls{llm} en el contexto de la preparación para esta certificación?
\end{itemize}

Estas preguntas permiten estructurar el proceso de evaluación del asistente virtual y sirven como base para el análisis de los resultados presentados en los capítulos posteriores, manteniendo una coherencia clara entre el problema identificado, la solución propuesta y las conclusiones del trabajo.

\section{Alcance y delimitaciones del trabajo}

El presente trabajo tiene como alcance el diseño, desarrollo e implementación de un asistente virtual basado en inteligencia artificial generativa, orientado a apoyar la preparación del examen de certificación \gls{pmp}. El asistente se concibe como una herramienta de apoyo educativo que complementa los métodos tradicionales de estudio, facilitando la comprensión de escenarios situacionales y el razonamiento alineado con el criterio del \gls{pmi}.

Desde el punto de vista técnico, el trabajo abarca el diseño de la arquitectura del sistema, la definición de los flujos de interacción, la implementación de un prototipo funcional y su evaluación en un contexto controlado de uso. El enfoque adoptado prioriza la viabilidad técnica y la utilidad pedagógica del asistente, por sobre el desarrollo de soluciones completamente integradas o de carácter comercial.

En cuanto al contenido, el asistente se orienta exclusivamente a temas vinculados con la preparación del examen \gls{pmp}, considerando los dominios, principios y enfoques promovidos por el \gls{pmi}. No se aborda la enseñanza general de la gestión de proyectos ni se pretende cubrir en profundidad otros marcos, metodologías o certificaciones relacionadas.

El trabajo no tiene como objetivo garantizar la aprobación del examen \gls{pmp} ni sustituir la formación formal, los cursos especializados o los materiales oficiales del \gls{pmi}. Asimismo, no se contempla el desarrollo de funcionalidades avanzadas tales como sistemas de recomendación personalizados basados en métricas de desempeño a largo plazo, análisis estadístico extensivo de resultados o integraciones completas con plataformas educativas externas.

Desde una perspectiva metodológica, la evaluación realizada se limita a un enfoque predominantemente cualitativo con un grupo reducido de participantes, por lo que los resultados obtenidos deben interpretarse dentro de este marco. Estas delimitaciones permiten acotar el alcance del trabajo, enfocando el análisis en los objetivos propuestos y asegurando la coherencia entre el problema planteado, la solución desarrollada y los resultados presentados.

\section{1.5 Organización del documento}
El presente informe se encuentra organizado en nueve capítulos, estructurados de manera progresiva con el objetivo de guiar al lector desde la contextualización del problema y el marco teórico, hasta el diseño, implementación, evaluación y conclusiones del asistente virtual desarrollado. Esta organización busca facilitar la comprensión integral del trabajo realizado, asegurando coherencia entre los objetivos planteados, la metodología adoptada y los resultados obtenidos.

A continuación, se presenta una breve descripción del contenido de cada uno de los capítulos que conforman el informe final:

\begin{itemize}
\item \textbf{Capítulo 1:} contextualiza el trabajo en el marco de la evolución reciente de la inteligencia artificial y, particularmente, de los \glspl{llm}. Se analiza su impacto en la educación y la formación profesional, así como la relevancia de la certificación \gls{pmp} en el ámbito laboral. Además, se expone el planteamiento del problema, se formulan la hipótesis y las preguntas de investigación, se delimita el alcance del trabajo y se describe la organización general del documento.
\item \textbf{Capítulo 2:} desarrolla los fundamentos conceptuales que sustentan la investigación. En este capítulo se abordan los principios de la gestión de proyectos y la certificación \gls{pmp}, los \glspl{llm} y sus capacidades, así como su aplicación en contextos educativos. Asimismo, se analizan técnicas de estudio relevantes para el aprendizaje autodirigido y se realiza un estudio comparativo de trabajos y herramientas relacionadas, destacando las diferencias con la propuesta desarrollada.
\item \textbf{Capítulo 3:} describe el enfoque metodológico adoptado para el desarrollo del proyecto. Se justifica la elección de un enfoque iterativo e incremental basado en el marco de trabajo Scrum, se detalla su adaptación al contexto del Trabajo Final y se presenta la planificación general del proyecto. Además, se describen en detalle los sprints ejecutados, vinculándolos con los objetivos específicos del trabajo.
\item \textbf{Capítulo 4:} se centra en el diseño conceptual y técnico del asistente virtual. Se definen los objetivos de diseño desde una perspectiva funcional, educativa y de experiencia de usuario, se adopta un enfoque de diseño centrado en el usuario y se describe la arquitectura general del sistema. Asimismo, se detallan los criterios y principios aplicados al diseño de la interfaz y al sistema de retroalimentación.
\item \textbf{Capítulo 5:} aborda los aspectos técnicos de la solución desarrollada. En este capítulo se justifica la selección del modelo de lenguaje utilizado, se describen las tecnologías empleadas en el backend y el frontend, y se explica la lógica educativa implementada en el asistente, incluyendo la generación de preguntas, la provisión de explicaciones y la adaptación al progreso del usuario.
\item \textbf{Capítulo 6:} presenta el diseño y la ejecución de los casos de estudio utilizados para evaluar el asistente virtual. Se describen los objetivos, los criterios de selección de los participantes y los instrumentos de evaluación. Posteriormente, se detallan los resultados obtenidos en cada caso, incluyendo la interacción de un aspirante a la certificación \gls{pmp} y la evaluación realizada por un profesional certificado.
\item \textbf{Capítulo 7:} expone los resultados obtenidos a partir de la aplicación de la metodología de evaluación definida. Se presentan y analizan tanto los resultados cuantitativos como cualitativos, utilizando métricas e instrumentos previamente establecidos, y se contrastan los hallazgos con los objetivos planteados y la hipótesis formulada.
\item \textbf{Capítulo 8:} ofrece una interpretación crítica de los resultados obtenidos. Se analizan los principales aportes del trabajo desde una perspectiva académica, tecnológica y educativa, se identifican las limitaciones del estudio y se reflexiona sobre las implicancias del uso de \glspl{llm} en el ámbito de la ingeniería en informática.
\item \textbf{Capítulo 9:} presenta las conclusiones generales del trabajo, evalúa el grado de cumplimiento de los objetivos propuestos y plantea posibles líneas de trabajo futuro. Estas incluyen la extensión funcional del asistente, su aplicación a otros dominios de certificación y la realización de evaluaciones a mayor escala.
\item \textbf{Anexos:} se incluyen materiales complementarios que enriquecen la documentación del trabajo, tales como detalles técnicos, diagramas, instrumentos de recolección de datos y capturas de pantalla de la solución desarrollada. Específicamente, el \textbf{Anexo A} detalla las Épicas e Historias de Usuario, el \textbf{Anexo B} presenta la planificación detallada de los sprints y el backlog de desarrollo, el \textbf{Anexo C} documenta la selección del modelo \gls{llm} y su comparación con alternativas, el \textbf{Anexo D} recopila los instrumentos de evaluación y protocolos de prueba utilizados, y el \textbf{Anexo E} proporciona el Manual de Uso del Asistente para el usuario final.
\end{itemize}
